\chapter{Acknowledgements}

% TODO: write acknowledgements
\begin{flushright}
\textit{Do you think I'd give up?\\
That this might've shook the love from me\\
\dots\\
How could you think, darlin', I'd scare so easily?\\
Now that it's done\\
There's not one thing that I would change\\
\dots\\
If someone asked me at the end\\
I'd tell them, ``Put me back in it''}

\vspace{0.5em}
--- \textbf{Hozier}, ``Francesca'' (2023)
\end{flushright}

\begin{comment}
    Thank the following people:
    Margo, Sasha, Thomas, Ivan. 
    Surbhi and Amee
    Milad, Alex, and Jack. 
    Sid, Shaurya, Joe, Rubia, Soo, Rut and Tanya, Satvik, Kjell, Joyce.
    Jimmy, Leeon, Neha, Vrushank
    Derek, Adam, Nacho, Felipe, Hayley, Anna, Giovanni, Anny G.?

    Family

    Lexapro, Vivance, and Xanax
    This research was supported by the Natural Sciences and Engineering Research
Council of Canada (NSERC)
\end{comment}

I never set out to do a PhD.
In fact, I was quite convinced that it wasn't the right path for me.
I had never studied computer science with any seriousness and was surprised when I got an offer to join the masters program at UBC.
There is a line from \textit{Everything Everywhere All at Once} that I think about often: ``We developed an algorithm that calculates which statistically improbable action will place you in a universe on the edge of your local cluster, which then slingshots you to the desired universe\dots\ The less sense it makes, the better.''
Saying yes to that offer made no sense at all.
I will forever be grateful to \textbf{Prof. Ivan Bechastnikh} for seeing potential in me and, by accepting me to the program, changing the trajectory of my life.

Another statistically improbable but fortunate event happened in 2018 when \textbf{Prof. Margo Seltzer} joined UBC.
I first met her when she visited what was then the NSS Lab, and I managed to grab a slot to discuss some provenance-adjacent work.
I was jittery before the meeting, but that melted away once I realized how kind and curious she was about my (somewhat flawed) project.
Naturally, I jumped at the opportunity to register for her graduate seminar that fall.
What started as a modest class project under her guidance gradually evolved and ultimately became the foundation of this dissertation.
I sought out Margo's mentorship, and working under her sheltering wing has been the defining privilege of my graduate career.
Her unwavering support, encouragement, and infectious enthusiasm for systems research not only carried this dissertation across the finish line, but shaped the kind of researcher and person I aspire to be.

I am also lucky to have \textbf{Prof. Sasha Fedorova} on my supervisory committee.
Her sharp insight and expansive experience with WiredTiger has been a huge help in getting me unstuck and finishing this project.
I would also like to thank \textbf{Prof. Thomas Pasquier} for his support and feedback that helped make this work better.

From the day I arrived in Vancouver, I have been fortunate to strike friendships that have survived much turmoil and many radio silences.
\textbf{Anna Scholtz, Hayley Guillou, and Giovanni Viviani}, you were there from the beginning, and your friendship has been one of the most grounding forces in my life here.
\textbf{\emph{Dr.} Surbhi Palande}, you have been a constant support, friend, and mentor.
There were \emph{many} moments when I found myself overcome with doubt and unable to see a way forward, and on those days your support dragged me back into the sunshine.
My fellow graduate journeymen: \textbf{Shaurya Patel, Sidhartha Agrawal, Joe Wonsil, Chris Chen, Soo Yi Lim, Tanya Prasad, Rut Vora, Satvik Vemuganti, and Xuechun Cao} have been a collective talisman that has kept me on the path to completion.
Having people who understand what this process demands has kept me sane and moving forward when I might otherwise have stalled.
Systopia lab has provided the safe and nurturing environment where I could thrive, and for that I thank each and every member.

Then there are the friends I made along the way.
\textbf{Derek Tan}, I thank you for the many fun hangs, museum tours, and gentle nudging to explore and develop hobbies that keep me grounded.
\textbf{Adam Geller, Felipe Gonzalez, and Ignacio Palma}, I thank you for showing me that there is light at the end of the tunnel and for reviving hope that things will turn out alright.
\textbf{Ruben Guerra}, you have been my closest confidant through some turbulent times.
You've always tried to make things better and quietly nudged me to believe in myself when I struggled to do so.

And then there are those who were there long before any of this began. 
\textbf{Jimmy Narang}, you were there the day this all started: you were the one who told me to apply to schools in Canada and UBC in particular.
None of this would have happened without you, and I am grateful for your constant presence and second opinions when I am spiralling.
\textbf{Leeon Passi, Neha Das, and Vrushank Kenkre}, you have been in my life for so long it's hard to imagine a time you weren't.
Some days, when I think of home, yours are the faces I see.
Bangalore was a special place for me, and none of it would have been half as meaningful without you lot.
Your friendship predates grad school, Vancouver, and any thought I had of pursuing this path, and it has been a constant I could always rely on.

Saved for last, but ironically my day ones, my family: thank you for your support and belief.
You have been patient for these many years I have been off chasing this sidequest, and I know \emph{Baba} would be proud that all three of his grandchildren are Doctors of some sort.

I would also like to thank Wellbutrin, Vyvanse, and Xanax, and the small army of doctors who have kept me alive and well.